% !TEX root = ../mainfile.tex
% !TEX encoding = UTF-8 Unicode
% !TEX spellcheck = en_US

\chapter{Discussion}
\label{ch:discussion}

This chapter draws the event-study and regression results together, relates them to the
literature, and sets out their implications and limits.

\section{Findings in Light of the Literature}
\label{sec:disc_literature}

The results speak directly to the six literatures set out in Chapter~\ref{ch:literature}.
The sharpest verdict is for the \emph{dollar-erosion} strand. \textcite{kamin_2025} and
\textcite{jiang_krishnamurthy_lustig_richmond_2026} read the 2025 tariff episode as evidence that the dollar has stopped
behaving as a safe haven, and \textcite{hassan_mertens_wang_zhang_2025} supply the theory: a trade war that
isolates U.S.\ goods markets erodes the currency's safety premium. The findings here both
confirm and bound that claim. They confirm it for the tariff shock, when the dollar depreciates
and its risk sensitivity collapses to zero, but they bound it by showing that the erosion is
not a general weakening: the very same dollar appreciates around the contemporaneous Hormuz
shock. What the dollar-erosion papers identify as a regime change is, on this evidence,
specific to U.S.-originated economic-policy shocks rather than a property of the post-2025
dollar as such.

For the \emph{safe-haven} tradition \parencite{ranaldo_soderlind_2010,habib_stracca_2012,debock_decarvalhofilho_2015} the evidence
is corroborative. The dollar performs as that literature predicts around both
military/oil-supply shocks, appreciating as risk rises, and the formal difference test cannot
separate the U.S.-implicated Hormuz response from the external Ukraine benchmark, so the
haven bid survives even U.S.\ involvement when the shock is a security one. The conditional,
crisis-contingent definition of \textcite{baur_lucey_2010} is thus vindicated for the dollar: its
haven status is real but state-dependent. That same framing organises the \emph{gold}
result. Gold hedges against the dollar rather than serving as a universal safe haven. It
bids when the tariff shock weakens the dollar and falls when the Hormuz shock strengthens it,
rising alongside the dollar only for the external Ukraine shock. This is exactly the state-dependence
Baur and Lucey's distinction implies.

The \emph{commodity-currency} literature \parencite{chen_rogoff_2003,kilian_2009} is borne out only in part. At the five- to twenty-day horizons the exporter--importer spread carries the sign of the crude move in the three main episodes, but the clearest separation occurs after the tariff shock, the Hormuz spread remains insignificant, and Ukraine provides weaker evidence only at twenty days (Section~\ref{sec:disc_oil}). Finally, the \emph{geopolitical-risk} literature
\parencite{caldara_iacoviello_2022} supplies both the measurement and a mechanism. The within-crisis
ceasefire reversal, which unwinds the dollar premium as the threat recedes, is the
threat-versus-realisation distinction Caldara and Iacoviello emphasise, mapped onto the
exchange rate.

\section{The Dollar's Safe Haven Is Conditional on Shock Type}
\label{sec:disc_hierarchy}

The central finding is that the dollar's safe-haven response is conditioned by the
\emph{type} of shock, not merely by who causes it. The dollar appreciated around both
military/oil-supply shocks, the external Ukraine invasion and the U.S.-involved Hormuz
crisis, and the formal difference test cannot distinguish the two responses at five and ten days, even though the United States was a belligerent in one and merely witnessed the other. Yet it depreciated around the
one purely economic-policy shock it created, the Liberation Day tariffs. Origin alone does
not explain the pattern. The nature of the shock does. This refines the ``dollar-erosion''
reading \parencite{kamin_2025,jiang_krishnamurthy_lustig_richmond_2026}: the safe-haven function is not weakening in
general but is specifically suspended when the United States is itself the source of an
economic-policy shock, while it survives intact, despite the subsequent reversion, for the
security shocks that the safe-haven tradition has always emphasised
\parencite{ranaldo_soderlind_2010,habib_stracca_2012}.

The regression sharpens this further. Estimated over the full sample, the dollar's
sensitivity to risk is significantly positive on normal days, reflecting the textbook safe-haven
reflex, but the interaction term shows it collapsing to essentially zero inside the tariff
window, while it remains intact during the Hormuz and Ukraine windows. The breakdown is thus
not merely a difference in the \emph{level} of the dollar across episodes but a change in its
very \emph{relationship with risk}: during the U.S.\ trade-policy shock the dollar stopped
trading like a safe haven, and only then.

Read against the hypotheses of Section~\ref{sec:hypotheses}, the evidence adjudicates
all three. \textbf{H1 is supported}: the dollar appreciated around both military/oil-supply
shocks, depreciated around the trade-policy shock, and the difference-in-CARs test separates
the tariff episode from both others while leaving Hormuz and Ukraine indistinguishable at the five- and ten-day horizons, with only
a difference of intensity on impact and a marginal gap at twenty days.
\textbf{H2 receives mixed evidence rather than robust support}: in the U.S.-originated tariff episode gold's ten-day cumulative abnormal return is positive in dollars ($+6.07\%$, significant at the 5\% level) but small and insignificant in euros ($+0.21\%$), and its five-day dollar-denominated estimate is negative and insignificant ($-1.54\%$), so the prediction is met at the ten-day dollar-denominated horizon and the result is sensitive to horizon and denomination. Gold's negative cumulative abnormal returns in the Hormuz crisis ($-22.30\%$ at twenty days, significant at the 1\% level) further weaken its interpretation as an origin-blind safe-haven benchmark (Section~\ref{sec:disc_oil}).
\textbf{H3 is partly supported}: at the five- to twenty-day horizons the exporter--importer spread carries the sign of the crude response in the three main episodes, negative under the tariff shock and positive under Hormuz and Ukraine, but the separation is significant at the 1\% level only for the tariff shock and, more weakly, at twenty days for Ukraine (Section~\ref{sec:disc_oil}).

\section{Gold and the Oil Channel}
\label{sec:disc_oil}

Gold does not behave as a universal safe haven, an asset that is uncorrelated or negatively correlated with risky assets in times of market stress in the sense of \textcite{baur_lucey_2010}, whatever the source of that stress. Instead it tracks the inverse of the dollar across the
two U.S.-linked shocks, rising when the tariff shock weakened the dollar, falling when the
Hormuz shock strengthened it, and bids alongside the dollar only for the external Ukraine
shock. Gold is thus better described as a hedge \emph{against} the dollar than as a
substitute for it, which reinforces, rather than competes with, the chapter's main message:
the dollar's own conditional move is the organising variable, and gold prices it from the
other side.

The oil channel sorts the currencies as the commodity-currency literature predicts only in part. The clearest exporter--importer separation came after the tariff announcement, when crude fell, with an exporter-minus-importer spread of $-2.80$ percentage points over days 0 to 10, significant at the 1\% level. Under Hormuz, when crude surged, the spread was positive but insignificant at five, ten and twenty days, and around the Ukraine invasion it reached $+4.38$ percentage points at twenty days, significant only at the 10\% level. The pattern is consistent with the finding of \textcite{chen_liu_wang_zhu_2016} that dollar exchange rates respond to demand-driven rather than supply-driven oil moves, although their monthly structural model and this daily event study identify different objects, and the overlap between the oil-importer and safe baskets (Section~\ref{sec:es_tariff}) prevents the spread from isolating oil exposure from safe-haven characteristics. For currencies such as the yen, safe-haven demand and higher oil-import costs may exert opposing pressures during an oil-supply shock.

\section{Policy Implications}
\label{sec:disc_policy}

Three implications follow. For reserve managers and central banks, the dollar's safe-haven
status cannot be treated as unconditional: hedging programmes calibrated on the assumption
that the dollar always rallies in risk-off episodes will misfire precisely when the shock
originates in U.S.\ economic policy. For portfolio managers, diversification should be
thought of across \emph{shock types}, not only across assets. Gold and competing-haven
currencies, not the dollar, carried the flight to quality in the tariff episode. For the
broader debate on the dollar's reserve role \parencite{obstfeld_zhou_2023,jiang_krishnamurthy_lustig_richmond_2026}, the evidence
is a caution against over-reading any single episode: the 2025 tariff shock did dent the
dollar, but the 2026 security shock showed the haven bid still firmly in place.

\section{Limitations}
\label{sec:disc_limitations}

Several caveats bound the results. Two Gulf events, the U.S.\ strikes of 25~May 2026 and the memorandum of understanding of 17~June 2026, are not analysed. The estimation window of the first contains 40 of its 120 trading days on or after the 2~March onset of the Hormuz crisis, so its benchmark would be contaminated, and the second lies too close to the sample end of 30~June 2026 for the longer windows (Section~\ref{sec:es_robust}). The regression's crisis windows are defined
around the events rather than by a structural break test, and the VIX--dollar link it
measures is a reduced-form association, not a causal channel. Gold's exceptional 2025--26
trend makes its constant-mean CARs
sensitive in level, so they are read for sign rather than magnitude. As a single-country
event study around a small number of recent, partly overlapping episodes, the design trades
breadth for the clean U.S.-versus-external contrast, and without forward/basis data beyond
the one-month tenor, covered-interest-parity deviations are left for future work.
